% =============================================================================
% Lecture 15b: Responsible Engineering, Part 2: Cost, Carbon, and Evidence
% Introduction to Machine Learning Systems (Volume I)
% =============================================================================
% Part 1 (Lecture 15) asked whether the system treats people correctly. Part 2
% asks what else it owes and whether it can prove it: the energy, carbon, and
% dollars of its whole life, and the duties it owes the people whose data it
% holds, each worked as a number the machine must compute, spend, and record.
%
% Build: make -C slides build-15b_responsible_engr   (figures: images/make_figs.py)
% Plan:  lecture_map.yml and storyboard.md in this folder.
% =============================================================================

\documentclass[aspectratio=169, 12pt]{beamer}
\usepackage{../../assets/beamerthememlsys}
\mlsyssetup{
  volume={Volume I},
  chapter={Chapter 15},
  chaptertitle={Responsible Engineering, Part 2},
  layer=applications,
  logo={../../assets/img/logo-mlsysbook.png},
  instlogo={../../assets/img/logo-harvard.png},
  dividerstack,
}
\usepackage[T1]{fontenc}
\usepackage[scaled=0.9]{helvet}
\usepackage{courier}
\renewcommand{\familydefault}{\sfdefault}
\usepackage{booktabs,amsmath}
\graphicspath{{images/}}
% Generated by calcs.py. Do not edit; edit calcs.py and rerun it.
\defval{nas_est_lb}{626{,}155}
\defval{car_lb}{126{,}000}
\defval{nas_est_t}{284}
\defval{car_t}{57}
\defval{nas_cars}{5}
\defval{nas_meas_t}{3.2}
\defval{nas_meas_t0}{3}
\defval{nas_err}{88}
\defval{nas_err_calc}{89}
\defval{proxy}{18.7}
\defval{tpu_perf_w}{3.3}
\defval{pue_g}{1.10}
\defval{pue_us}{1.58}
\defval{pue_ratio}{1.44}
\defval{grid_g}{0.431}
\defval{grid_us}{0.423}
\defval{factors}{88}
\defval{proxy_batch}{4{,}096}
\defval{full_batch}{32{,}768}
\defval{nas_iowa_t}{0.6}
\defval{nas_days}{6.8}
\defval{nas_chips}{200}
\defval{nas_w}{208}
\defval{nas_kwh}{7{,}468}
\defval{nas_calc_t}{3.22}
\defval{meena_t}{48.5}
\defval{payback}{15}
\defval{payback_1}{15.2}
\defval{hook_cost}{17.7}
\defval{hook_health}{46.5}
\defval{fb_gap0}{3}
\defval{fb_alpha}{1.10}
\defval{fb_alpha_1}{1.1}
\defval{fb_gate}{5}
\defval{fb_cross}{5.4}
\defval{fb_fail_q}{6}
\defval{fb_gap_fail}{5.3}
\defval{credit_all_n}{12{,}000}
\defval{credit_appr_n}{6{,}300}
\defval{gem_wh}{0.24}
\defval{gem_narrow_wh}{0.10}
\defval{gem_pue}{1.09}
\defval{gem_s2}{0.023}
\defval{gem_s13}{0.010}
\defval{gem_intensity}{0.096}
\defval{gem_total}{0.03}
\defval{gem_energy_x}{33}
\defval{gem_carbon_x}{44}
\defval{gem_model_x}{23}
\defval{gem_util_x}{1.4}
\defval{ff_base_t}{0.1357}
\defval{ff_base_kwh}{316}
\defval{ff_best_t}{0.0024}
\defval{ff_best_kwh}{30}
\defval{ff_printed}{57}
\defval{ff_ops}{1.6}
\defval{ff_model}{1.3}
\defval{ff_processor}{7.4}
\defval{ff_datacenter}{10.5}
\defval{ff_grid}{56.5}
\defval{ff_proc_step}{5.7}
\defval{ff_pue_a}{1.59}
\defval{ff_pue_b}{1.11}
\defval{ff_pue_step}{1.4}
\defval{ff_grid_a}{0.429}
\defval{ff_grid_b}{0.080}
\defval{ff_grid_step}{5.4}
\defval{ff_iowa_local}{0.478}
\defval{ff_cfe_pct}{78}
\defval{gpt3_mwh}{1{,}287}
\defval{gpt3_kwh}{1{,}287{,}000}
\defval{gpt3_pue}{1.1}
\defval{gpt3_grid}{0.429}
\defval{gpt3_t}{552}
\defval{car_year_t}{4.6}
\defval{gpt3_cars}{120}
\defval{bloom_mwh}{433}
\defval{bloom_gpu_kwh}{433{,}195}
\defval{bloom_pue}{1.2}
\defval{bloom_grid}{0.057}
\defval{bloom_grid_g}{57}
\defval{bloom_kwh}{520{,}000}
\defval{bloom_kmwh}{520}
\defval{bloom_t}{29.6}
\defval{energy_x}{2.48}
\defval{energy_x3}{2.476}
\defval{grid_x}{7.53}
\defval{grid_x3}{7.526}
\defval{carbon_x}{18.6}
\defval{naive_energy_x}{3.0}
\defval{naive_carbon_x}{22}
\defval{bloom_t4}{30}
\defval{llm_gap_t4}{18}
\defval{llm_gpt3}{552}
\defval{llm_gopher}{380}
\defval{llm_opt}{76.3}
\defval{llm_bloom}{30}
\defval{llm_span}{18}
\defval{pt_min_t}{4.3}
\defval{pt_max_t}{552}
\defval{pt_min_grid}{0.177}
\defval{pt_max_grid}{0.545}
\defval{b_dyn}{24.7}
\defval{b_idle}{14.6}
\defval{b_emb}{11.2}
\defval{b_total}{50.5}
\defval{b_dyn_share}{48.9}
\defval{b_dyn_share0}{49}
\defval{b_off_kwh}{27}
\defval{b_idle_kwh}{64}
\defval{b_train_kwh}{109}
\defval{b_code_pct}{41}
\defval{b_idle_total_kwh}{256{,}646}
\defval{b_servers_t}{7.57}
\defval{b_gpus_t}{3.64}
\defval{server_kg}{2{,}500}
\defval{gpu_kg}{150}
\defval{life_years}{6}
\defval{use_pct}{85}
\defval{server_kg_h}{0.056}
\defval{gpu_kg_h}{0.0034}
\defval{card_g}{24{,}700{,}000}
\defval{card_t}{24.7}
\defval{card_share}{49}
\defval{users_m}{10}
\defval{items}{20}
\defval{inf_day_m}{200}
\defval{service_ms}{10}
\defval{service_s}{0.01}
\defval{gpu_s_day_m}{2}
\defval{gpu_h_day}{556}
\defval{gpu_h_day_1}{555.6}
\defval{gpus}{23.1}
\defval{gpus_0}{23}
\defval{servers}{2.9}
\defval{gpus_per_server}{8}
\defval{servers_w_words}{three}
\defval{inf_rate}{2.50}
\defval{train_rate}{4}
\defval{serve_year}{506{,}944}
\defval{serve_year_k}{507}
\defval{cycle}{3{,}200}
\defval{cycle_prep}{400}
\defval{cycle_search}{2{,}000}
\defval{cycle_final}{800}
\defval{cycles}{twelve}
\defval{years}{3}
\defval{days}{1{,}095}
\defval{train_3y_k}{38}
\defval{train_3y_k1}{38.4}
\defval{inf_3y_m}{1.52}
\defval{ops_year_k}{170}
\defval{ops_3y_k}{510}
\defval{tco_m}{2.07}
\defval{train_share}{1.9}
\defval{inf_share}{73.5}
\defval{cost_x}{40}
\defval{gpu_w}{400}
\defval{gpu_kw}{0.4}
\defval{fleet_grid}{0.4}
\defval{kg_gpu_h}{0.16}
\defval{inf_gpu_h}{608{,}333}
\defval{inf_t}{97.3}
\defval{inf_t0}{97}
\defval{train_t}{1.5}
\defval{carbon_x_dlrm}{63}
\defval{inf_t_year}{32}
\defval{serve_cut_pct}{20}
\defval{serve_cut_frac}{0.2}
\defval{train_cut_frac}{0.5}
\defval{save_serve_k}{304}
\defval{save_serve_t}{19}
\defval{save_train_k}{19}
\defval{save_train_t}{0.8}
\defval{save_serve_pct}{15}
\defval{save_train_pct}{1}
\defval{halvings}{16}
\defval{final_half}{400}
\defval{frontier_growth}{2.9}
\defval{priciest_m}{388}
\defval{api_req}{230{,}768}
\defval{api_kwh}{914}
\defval{api_gpus}{16}
\defval{api_days}{18}
\defval{api_wh}{3.96}
\defval{api_idle_kwh}{0.28}
\defval{api_loaded_pct}{75}
\defval{g_inf_share}{60}
\defval{g_train_share}{40}
\defval{g_ml_lo}{10}
\defval{g_ml_hi}{15}
\defval{m_exp}{10}
\defval{m_train}{20}
\defval{m_inf}{70}
\defval{m_emb}{30}
\defval{op_dirty_g}{160}
\defval{op_clean_g}{23}
\defval{emb_g}{10.4}
\defval{gpu_g_h}{3.4}
\defval{server_g_h}{56}
\defval{emb_dirty_pct}{6}
\defval{emb_clean_pct}{31}
\defval{dlrm_emb_t}{6.3}
\defval{dlrm_emb_t_year}{2}
\defval{dpc_fb}{210}
\defval{dpc_ig}{180}
\defval{dpc_total}{390}
\defval{decisions_m}{1}
\defval{deny_pct}{47.5}
\defval{deny_frac}{0.475}
\defval{samples}{1{,}000}
\defval{features}{20}
\defval{notices}{475{,}000}
\defval{sampled_m}{475}
\defval{sampled_x}{475}
\defval{exact_per}{1{,}048{,}576}
\defval{exact_mant}{5}
\defval{exact_exp}{11}
\defval{all_decisions_b}{1}
\defval{appeal_pct}{0.1}
\defval{appeals}{1{,}000}
\defval{case_min}{6}
\defval{review_h}{100}
\defval{reviewers}{12.5}
\defval{shift_h}{8}
\defval{cigna_cases}{300{,}000}
\defval{cigna_s}{1.2}
\defval{cigna_h}{100}
\defval{cigna_months}{two}
\defval{credit_per_h}{10}
\defval{cigna_per_h}{3{,}000}
\defval{carlini_mem}{604}
\defval{carlini_cand}{1{,}800}
\defval{carlini_names}{46}
\defval{carlini_contact}{32}
\defval{cycle_days}{91}
\defval{cycle_days_r}{90}
\defval{sisa_purchase}{4.63}
\defval{sisa_svhn}{2.45}
\defval{sisa_imagenet}{1.36}
\defval{eps}{8}
\defval{e_eps}{3{,}000}
\defval{kws_params}{200{,}000}
\defval{kws_params_k}{200}
\defval{fp32_b}{4}
\defval{kws_grad_kb}{800}
\defval{batch}{256}
\defval{kws_batch_mb}{204.8}
\defval{gpt2_params_b}{1.5}
\defval{gpt2_grad_gb}{6}
\defval{gpt2_batch_tb}{1.5}
\defval{gpt2_batch_tb3}{1.536}
\defval{h100_gb}{80}
\defval{cifar_dp}{81.4}
\defval{cifar_prev}{71.7}
\defval{cifar_gain}{9.7}
\defval{cifar_np}{95}
\defval{cifar_gap}{13}
\defval{im_dp}{86.7}
\defval{im_same}{88.5}
\defval{im_gap_same}{1.8}
\defval{im_best}{91.0}
\defval{im_gap_best}{4.3}
\defval{im_eps05}{83.8}
\defval{audio_khz}{16}
\defval{audio_bits}{16}
\defval{audio_bps}{32{,}000}
\defval{audio_mb_h}{115.2}
\defval{update_mb}{0.8}
\defval{audio_x}{144}
\defval{rec_b}{256}
\defval{ret_months}{25}
\defval{ret_days}{760}
\defval{log_mb_day}{256}
\defval{log_gb}{195}
\defval{dlrm_req_m}{200}
\defval{dlrm_days}{365}
\defval{dlrm_manifest_tb}{18.7}
\defval{dense}{13}
\defval{sparse}{26}
\defval{emb_dim}{128}
\defval{input_b}{13{,}364}
\defval{dlrm_input_tb}{976}
\defval{payload_x}{52}
\defval{cc_wh}{0.186}
\defval{cc_mg}{68.6}
\defval{cc_intensity}{0.369}
\defval{cc_cpu_pct}{73}

\setcounter{mlsystotalsections}{4}

\title{Responsible Engineering, Part 2}
\author{Vijay Janapa Reddi}
\institute{Harvard University}
\date{}

% A large answer choice for predict and concept-question slides
\newcommand{\choice}[2]{\textbf{\textcolor{mlsysaccent}{#1}}\enspace #2}
\newcommand{\rchoice}[2]{\colorbox{mlsysaccent!15}{\choice{#1}{#2}}}
% A small gray tag that marks a constructed, assumed, or derived value
\newcommand{\smalltag}[1]{{\scriptsize\textcolor{midgray}{\textit{#1}}}}
% A big number with a caption under it
% A takeaway line with its text-free thumbnail
\newcommand{\thumbline}[2]{%
  \begin{minipage}[c]{2.6cm}\centering\includegraphics[width=2.5cm,height=0.85cm,keepaspectratio]{#1}\end{minipage}%
  \hspace{0.3cm}\begin{minipage}[c]{0.74\textwidth}\raggedright\small #2\end{minipage}}
\newcommand{\bignum}[2]{{\fontsize{34}{38}\selectfont\bfseries\textcolor{mlsysaccent}{#1}}\\[0.1cm]{\small #2}}


\begin{document}

% =============================================================================
% TITLE, HOOK, RECALL
% =============================================================================
\mlsysnote{[0.5 min] Part 2 of Responsible Engineering. Part 1 asked whether the system treats people correctly. Today asks what else the system owes, its energy, carbon, and money over its whole life and its duties to the people whose data it holds, and whether it can prove each one. Every number today is worked step by step, with the same rigor as a latency budget. Four running cases: large training runs (measured), Google's Gemini prompt (measured), the book's DLRM ranking service (a scenario), and the credit service and voice assistant from Part 1. The notes total about 71 minutes; cutting the frames tagged [OPTIONAL] leaves a complete lecture of about 49.5 minutes. FLEX: [CORE].}
\mlsystitle{Responsible Engineering, Part~2}{What a system consumes over its life, and what it owes the people whose data it holds, are numbers the machine must compute and record.}{cover_responsible_engr_blueprint.png}

\begin{frame}{Predict: how much did one architecture search emit?}
  \mlsysactivity{Predict}{3 min}
  \centering
  \mlsyssteps[width=0.92\textwidth,height=3.9cm,keepaspectratio]{hook_nas_step}{3}\\[0.1cm]
  \only<2|handout:0>{\small The team that ran the search measured it in 2021. The real number was\\[0.05cm]
    \normalsize\choice{A}{600 t}\qquad\choice{B}{\val{nas_est_t} t}\qquad\choice{C}{30 t}\qquad\choice{D}{\val{nas_meas_t0} t}}%
  \only<3|handout:0>{\normalsize\choice{A}{600 t}\qquad\choice{B}{\val{nas_est_t} t}\qquad\choice{C}{30 t}\qquad\rchoice{D}{\val{nas_meas_t0} t}}%
  \only<4-|handout:1>{\small Why: a small proxy task (\val{proxy}$\times$); TPU v2 in place of P100 (\val{tpu_perf_w}$\times$ per watt);\\PUE \val{pue_g} against \val{pue_us}. The grid was about the same.}
  \mlsyscredit{Strubell et al., ACL 2019, Table 1 (\val{nas_est_lb} lb; car lifetime \val{car_lb} lb). Patterson et al., ``Carbon Emissions and Large Neural Network Training,'' 2021, Table 4, \S4.1, and Appendices C and D.}
  \note{[4 min] Step 1: in 2019 Strubell and colleagues estimated that the architecture search that found the Evolved Transformer emitted \val{nas_est_lb} pounds of CO2e (CO2 equivalent: all greenhouse gases in CO2 terms), \val{nas_est_t} tonnes, beside one car's lifetime at \val{car_t} tonnes. Step 2: ENGAGE. The team that ran the search measured it in 2021. Step 3: D, \val{nas_meas_t} tonnes, \val{nas_err} times lower (Patterson et al. 2021, Table 4 and section 4.1). Note the log axis. WARN: students who chose B treat a published estimate as a measurement. Step 4: three causes (the paper's Appendices C and D). The search used a small proxy task, Transformer (Base) at batch \val{proxy_batch} in place of (Big) at \val{full_batch}, which accounts for \val{proxy} times, the authors' error for an average organization. It ran on TPU v2 chips in place of P100 GPUs, \val{tpu_perf_w} times better performance per watt, from measured power. PUE, power usage effectiveness (building energy over computer energy), was \val{pue_g} against the 2018 US average of \val{pue_us}. Grid intensity, CO2e per kWh delivered, was about the same: \val{grid_g} kg against \val{grid_us}. \val{proxy} times \val{tpu_perf_w} times \val{pue_ratio} is \val{factors} (calcs.py). The same search in Google's Iowa datacenter, counted after its carbon-free energy purchases, would have emitted \val{nas_iowa_t} tonnes. No one had published which task, processor, or datacenter. FLEX: [CORE].}
\end{frame}

\begin{frame}{Recall: what did Part 1 measure?}
  \mlsysactivity{Retrieval}{1 min}
  \begin{columns}[c]
    \begin{column}{0.56\textwidth}
      \small 90 seconds.\\[0.25cm]
      \begin{enumerate}\setlength{\itemsep}{0.15cm}
        \item Ranked by health, what share of auto-enrolled patients would be Black?
        \item Name the three fairness criteria and the rate each equalizes.
        \item A \val{fb_gap0}-point gap grows by a factor of \val{fb_alpha} each quarter. In which quarter does it cross \val{fb_gate}?
      \end{enumerate}
    \end{column}
    \begin{column}{0.42\textwidth}
      \uncover<2->{\renewcommand{\arraystretch}{1.5}\small
      \begin{tabular}{@{}rl@{}}
        1. & \textbf{\val{hook_health}\%} against \val{hook_cost}\% \\
        2. & approval rate; TPR; TPR and FPR \\
        3. & \textbf{quarter \val{fb_fail_q}}, at \val{fb_gap_fail} points
      \end{tabular}}
    \end{column}
  \end{columns}
  \mlsyscredit{Obermeyer et al., Science 2019, Fig. 1B; Responsible Engineering, \S15.3.2.1 and Napkin Math 15.3.}
  \note{[1.5 min] Silent writing, then step 2. Answers: \val{hook_health} percent against \val{hook_cost} percent under the cost label. Demographic parity equalizes approval rates, equal opportunity the true positive rate, equalized odds the true and false positive rates. \val{fb_gap0} times \val{fb_alpha_1} to the k exceeds \val{fb_gate} at k above \val{fb_cross}, so quarter \val{fb_fail_q}, at \val{fb_gap_fail} points. Part 1 asked whether the system treats people correctly. Today asks what else it owes. A student who missed Part 1 needs only the credit service: a constructed lender with \val{credit_all_n} applicants in the worked example, and \val{decisions_m} million decisions a day in today's governance scenario; both are reintroduced when used. FLEX: [CORE].}
\end{frame}

\begin{frame}{What else does a system owe, and can it prove it?}
  \begin{columns}[c]
    \begin{column}{0.76\textwidth}
      \normalsize By the end of today, you will be able to:\\[0.1cm]
      \small
      \begin{itemize}\small\setlength{\itemsep}{0.1cm}
        \item Convert measured energy to carbon with PUE and grid intensity, and add idle and embodied terms.
        \item Find the term that dominates a service's lifetime cost and carbon.
        \item Turn a legal duty into a capability and estimate its machine cost.
        \item Trace an erasure request, and say what weights, unlearning, and privacy cost.
      \end{itemize}
    \end{column}
    \begin{column}{0.22\textwidth}
      \centering\mlsyskitfig[width=\textwidth,height=3.2cm,keepaspectratio]{stack_applications}
    \end{column}
  \end{columns}
  \vspace{0.2cm}
  \threadstrip{0}{1}{0}{0}{0}
  \note{[1.5 min] The title is the driving question. A responsible system answers for more than its decisions: what energy, carbon, and money it consumes over its life, what it owes the people whose data it holds, and what the machine must compute and keep to prove both. Read the four outcomes. The stack shows where we are, the applications layer, where the costs of every layer below are added up and owed to someone. LINK: Lecture 2 named total cost of ownership when comparing cloud and edge; today defines its lifecycle split. The strip fills the fast-enough half of the dual mandate, read as a budget of energy, carbon, and money. FLEX: [CORE].}
\end{frame}

% =============================================================================
% SECTION: JOULES TO TONNES
% =============================================================================
\mlsysnote{[0.5 min] Start with the law that converts energy to carbon, apply it to one computation, then let each factor move. FLEX: [CORE].}
\mlsysdivider{What sets the carbon of one computation?}{Responsible Engineering, \S15.4.1 and \S15.4.3}
\section{Joules to Tonnes}

\mlsysnote{[0.5 min] The lecture's one law. Read it with units: energy in kWh measured at the IT equipment, PUE a ratio of facility energy to IT energy, grid intensity in kg CO2e per kWh. Each factor is a lever a team controls: the model and processor set the energy, the datacenter sets PUE, the region and the hour set the intensity. Manufacturing is a separate term added beside it. The next slide applies the law to one measured prompt. FLEX: [CORE].}
\mlsysfocus{Operational carbon}{%
  {\Large Carbon $=$ Energy $\times$ PUE $\times$ Grid intensity}\\[0.5cm]
  {\small\textcolor{white!75}{kWh at the IT equipment \quad facility / IT energy \quad kg CO$_2$e per kWh}}%
}

\begin{frame}{One Gemini prompt: \val{gem_wh} Wh at the full boundary, \val{gem_total} g of CO\texorpdfstring{\textsubscript{2}}{2}e}
  \centering
  \mlsyssteps[width=0.90\textwidth,height=4.0cm,keepaspectratio]{gemini_step}{2}\\[0.05cm]
  \uncover<2->{\small $\val{gem_wh}\text{ Wh} \times \val{gem_intensity}\text{ kg/kWh} = \val{gem_s2}$ g\quad\smalltag{(intensity derived from the paper's two numbers)}\\[0.02cm]
  {\footnotesize\textcolor{midgray}{Scope 2: purchased electricity, market-based. Scopes 1 and 3: direct emissions and manufacturing.}}}\\[0.05cm]
  \uncover<3->{\small\textcolor{mlsysaccent}{\textbf{May 2024 to May 2025: \val{gem_energy_x}$\times$ less energy, \val{gem_carbon_x}$\times$ less carbon per prompt.}}}
  \mlsyscredit{Elsworth et al., ``Measuring the Environmental Impact of Delivering AI at Google Scale,'' 2025, Tables 1 and 2 and \S4.3. Median Gemini Apps text prompt, May 2025.}
  \note{[2 min] A measured production number, recalled from Lecture 12. Step 1: Google's median Gemini text prompt in May 2025 used \val{gem_wh} Wh at the full boundary against \val{gem_narrow_wh} at the accelerators alone; the facility share is a PUE of \val{gem_pue}, so \val{gem_wh} Wh is already energy times PUE. Today adds carbon. Step 2: carbon. The electricity term is \val{gem_s2} g; dividing by \val{gem_wh} Wh gives \val{gem_intensity} kg per kWh, the grid intensity Google reports after its clean-energy purchases. Read the gray line: this is market-based Scope 2; location-based accounting would use the local grid's average and give a larger number. Scopes 1 and 3 add \val{gem_s13} g, mostly the manufacturing of the hardware, so the total is \val{gem_total} g. Step 3: in one year the same prompt used \val{gem_energy_x} times less energy and emitted \val{gem_carbon_x} times less carbon, \val{gem_model_x} times from model changes and \val{gem_util_x} times from keeping machines busier. WARN: \val{gem_wh} Wh is a median over prompts on one service; it is not comparable with a number measured at a narrower boundary. FLEX: [OPTIONAL]; cutting it loses the per-prompt example.}
\end{frame}

\begin{frame}{Four choices cut one training job's carbon \val{ff_printed}-fold}
  \begin{columns}[T]
    \begin{column}{0.55\textwidth}
      \centering
      \mlsyssteps[width=\textwidth,height=5.0cm,keepaspectratio]{factors_step}{4}\\[0.05cm]
      {\footnotesize \val{ff_base_t} t on P100 in an average US datacenter $\rightarrow$ \val{ff_best_t} t}
    \end{column}
    \begin{column}{0.40\textwidth}
      \centering
      \uncover<3->{\safeimg[width=\textwidth,height=2.15cm,keepaspectratio]{google-iowa-council-bluffs.jpg}\\[0.02cm]
      {\scriptsize Google, Council Bluffs, Iowa: PUE \val{ff_pue_b}}\\[0.1cm]}
      \uncover<4->{\safeimg[width=\textwidth,height=2.25cm,keepaspectratio]{owid-carbon-intensity-europe-2022-map.png}\\[0.02cm]
      {\scriptsize Grid intensity, Europe 2022: France pale, Poland dark}}
    \end{column}
  \end{columns}
  \mlsyscredit{Patterson et al. 2021, Table 1 and Fig.\ 1 (as printed). Photo: Chad Davis, CC BY 2.0, Wikimedia Commons. Map: Our World in Data (Ember), lifecycle intensity 2022, CC BY 4.0.}
  \note{[2.5 min] [CORE] Worked example: the same accuracy, four ways. The baseline trains Transformer (Big) on NVIDIA P100 GPUs in an average US datacenter: \val{ff_base_kwh} kWh and \val{ff_base_t} tonnes. Step 1: model. The Evolved Transformer (Medium) reaches the same accuracy with \val{ff_ops} times fewer operations: \val{ff_model} times. Step 2: processor. TPU v2: \val{ff_processor} times cumulative, about \val{ff_proc_step} from the processor. Step 3: datacenter. PUE from \val{ff_pue_a} to \val{ff_pue_b} is \val{ff_pue_step} times; cumulative \val{ff_datacenter}. The photo: Google's Council Bluffs, Iowa datacenter, the paper's best case; PUE belongs to a building. Step 4: grid. \val{ff_grid_a} kg per kWh to \val{ff_grid_b} is \val{ff_grid_step} times; cumulative \val{ff_grid}, which the paper rounds to \val{ff_printed}: \val{ff_best_kwh} kWh and \val{ff_best_t} tonnes. WARN: the grid step is market-based. Iowa's local grid was \val{ff_iowa_local} kg per kWh, higher than the US average; the \val{ff_grid_b} is net of Google's carbon-free energy purchases, \val{ff_cfe_pct} percent hour by hour in 2019. Location-based accounting would not show that step. The map (Our World in Data, 2022) shows why region matters. France, mostly nuclear, is pale; Poland, mostly coal, dark. Its basis is lifecycle, so its values are not the operational \val{ff_grid_a} and \val{ff_grid_b}; read it for the spread. ENGAGE: which factor can a small team control? Expected: model and region, sometimes processor; rarely PUE. FLEX: [CORE].}
\end{frame}

\begin{frame}{Estimate: how much less carbon did BLOOM emit?}
  \mlsysactivity{Estimate}{3 min}
  \begin{columns}[T]
    \begin{column}{0.46\textwidth}
      \begin{mlsyscard}{mlsysaccent}
        \textbf{Given, two 175B-scale models}\\[0.1cm]
        GPT-3: \val{gpt3_mwh} MWh, PUE \val{gpt3_pue} included\\US grid \val{gpt3_grid} kg/kWh\\[0.15cm]
        BLOOM: \val{bloom_mwh} MWh at the GPUs\\PUE \val{bloom_pue}; French grid \val{bloom_grid} kg/kWh
      \end{mlsyscard}
    \end{column}
    \begin{column}{0.51\textwidth}
      \small
      \uncover<2->{GPT-3: $\val{gpt3_kwh} \times \val{gpt3_grid} = \mathbf{\val{gpt3_t}}$ t\\[0.1cm]
                   BLOOM: $\val{bloom_gpu_kwh} \times \val{bloom_pue} = \val{bloom_kwh}$ kWh\\
                   \quad $\val{bloom_kwh} \times \val{bloom_grid} = \mathbf{\val{bloom_t}}$ t\\[0.25cm]}
      \uncover<3->{Energy: $\val{gpt3_mwh} / \val{bloom_kmwh} = \val{energy_x}\times$\\[0.05cm]
                   Grid: $\val{gpt3_grid} / \val{bloom_grid} = \val{grid_x}\times$\\[0.05cm]
                   Carbon: $\val{energy_x3} \times \val{grid_x3} = \textbf{\textcolor{mlsysaccent}{\val{carbon_x}}}\times$}
    \end{column}
  \end{columns}
  \mlsyscredit{Responsible Engineering, Napkin Math 15.5 (GPT-3, after Patterson et al. 2021, Table 4, whose energy includes PUE); Luccioni et al., ``Estimating the Carbon Footprint of BLOOM,'' JMLR 2023, \S4.2 and Table 4.}
  \note{[2.5 min] ENGAGE: one minute with a neighbor; ask for a single number. GPT-3's energy is the number Lecture 0 gave. Step 2: first put both energies at the same boundary. GPT-3's \val{gpt3_mwh} MWh already includes its datacenter's PUE of \val{gpt3_pue}: Patterson and colleagues multiply chip-hours by measured power and by PUE. Times \val{gpt3_grid}, it is \val{gpt3_t} tonnes, \val{gpt3_cars} passenger-car years at \val{car_year_t} tonnes each (Napkin Math 15.5). BLOOM's \val{bloom_gpu_kwh} kWh is the GPUs' energy alone, so multiply by its datacenter's PUE of \val{bloom_pue}: about \val{bloom_kwh} kWh. This is the law's PUE factor used for the first time. On France's mostly nuclear grid at \val{bloom_grid_g} g per kWh, that is \val{bloom_t} tonnes. Step 3: the factors multiply. \val{energy_x} times less energy and \val{grid_x} times lower intensity give \val{carbon_x} times less carbon; Luccioni's table rounds BLOOM to \val{bloom_t4} tonnes, which gives the next slide's \val{llm_gap_t4}. WARN: students who answer 2.5 forgot the grid; the grid moved the result more than the energy did. WARN: dividing \val{gpt3_mwh} by \val{bloom_mwh} directly gives \val{naive_energy_x} and \val{naive_carbon_x} times, because it compares a building's energy with its computers' energy. Check the boundary before dividing. FLEX: [CORE].}
\end{frame}

\begin{frame}{Four large language models span \val{llm_span} times in training carbon}
  \centering
  \mlsyssteps[width=0.92\textwidth,height=4.8cm,keepaspectratio]{llm4_step}{2}\\[0.05cm]
  {\footnotesize * inferred by Luccioni et al.; the model's authors did not state it}
  \mlsyscredit{Luccioni, Viguier, Ligozat, JMLR 2023, Table 4 (CO\textsubscript{2}e $\times$ PUE). Starred values are in italics there: inferred by the authors from the models' papers.}
  \note{[1.5 min] Step 1: the two models from the estimate, as Luccioni's table reports them with each datacenter's PUE: GPT-3 \val{llm_gpt3} tonnes, BLOOM \val{llm_bloom}. Step 2: Gopher, \val{llm_gopher} tonnes, and OPT, \val{llm_opt}. Models of 175 to 280 billion parameters span \val{llm_span} times, and the label beside each bar says why: energy and grid. WARN: the stars matter. Luccioni and colleagues had to infer Gopher's energy and most of OPT's numbers, because those papers did not report them. A comparison table is only as good as what each team recorded, which is the hook's point again. For reference, Patterson's 2021 Table 4 gives measured net emissions for five Google and OpenAI models from \val{pt_min_t} tonnes (GShard-600B) to \val{pt_max_t} (GPT-3), with grid intensities from \val{pt_min_grid} to \val{pt_max_grid} kg per kWh. FLEX: [OPTIONAL]; the estimate carries the lesson.}
\end{frame}

\begin{frame}{Half of BLOOM's footprint was not the training computation}
  \centering
  \mlsyssteps[width=0.92\textwidth,height=4.0cm,keepaspectratio]{bloom_life_step}{3}\\[0.25cm]
  \uncover<2->{\small Idle: the cluster drew \val{b_idle_kwh} kWh per interval with nothing running, \val{b_train_kwh} while training.}\\[0.05cm]
  \uncover<3->{\small Embodied: servers \val{b_servers_t} t and GPUs \val{b_gpus_t} t, spread over a \val{life_years}-year life.}
  \mlsyscredit{Luccioni et al., JMLR 2023, Table 3 (life-cycle total), Table 2 and \S4.3 (idle), \S4.1 (embodied).}
  \note{[2.5 min] [CORE] The only full life-cycle account of a large model's training. Step 1: dynamic consumption, the energy the GPUs drew while running BLOOM's code, \val{b_dyn} tonnes, \val{b_dyn_share} percent. This is the number most papers report, and the estimate's starting energy before PUE. Step 2: idle consumption, \val{b_idle} tonnes. The team measured the cluster in three states: with compute nodes off, \val{b_off_kwh} kWh per measurement interval for network, storage, and cooling; with nodes on but idle, \val{b_idle_kwh}; while training, \val{b_train_kwh}. Of the \val{b_train_kwh} kWh drawn per interval while training, \val{b_idle_kwh} were drawn with the nodes on and nothing running, so the code added only about \val{b_code_pct} percent of training-time power (derived in calcs.py). Over the run, the idle draw added \val{b_idle_total_kwh} kWh. Step 3: embodied emissions, \val{b_emb} tonnes: the carbon of manufacturing, spread over the hardware's life. A server is about \val{server_kg} kg CO2e to make and a GPU at least \val{gpu_kg} kg; replaced every \val{life_years} years at \val{use_pct} percent use, that is \val{server_kg_h} kg per server-hour and \val{gpu_kg_h} kg per GPU-hour (calcs.py). Total \val{b_total} tonnes, twice the dynamic number. ENGAGE: is idle power a PUE term? Expected: no; PUE counts the building's overhead, idle counts IT equipment powered but not computing. WARN: embodied carbon is a separate term outside PUE. FLEX: [CORE].}
\end{frame}

\begin{frame}[fragile]{BLOOM's model card reports the dynamic term alone}
  \mlsysactivity{Find}{1.5 min}
  \begin{columns}[T]
    \begin{column}{0.60\textwidth}
      {\footnotesize\textbf{README.md}, metadata block (bigscience/bloom)}
\vspace{-0.1cm}
\begin{semiverbatim}\scriptsize
co2_eq_emissions:
\textcolor{RedLine}{  emissions: 24700000}
  source: "Estimating the Carbon Footprint of
    BLOOM, a 176B Parameter Language Model. ..."
  training_type: "pre-training"
\textcolor{BlueLine}{  geographical_location: "Orsay, France"}
\textcolor{BlueLine}{  hardware_used: "384 A100 80GB GPUs"}
\end{semiverbatim}
      {\footnotesize\textbf{Environmental Impact} section, same file}
\vspace{-0.1cm}
\begin{semiverbatim}\scriptsize
**Estimated carbon emissions:**  *(Forthcoming.)*
**Estimated electricity usage:** *(Forthcoming.)*
\end{semiverbatim}
    \end{column}
    \begin{column}{0.38\textwidth}
      \small Which line of a carbon record is here, and which is missing?\\[0.25cm]
      \uncover<2->{\textcolor{RedLine}{\textbf{\val{card_g} g $=$ \val{card_t} t}}: the dynamic term, \val{card_share}\% of \val{b_total} t\\[0.15cm]
      \textcolor{BlueLine}{\textbf{Present}}: hardware, location\\[0.1cm]
      \textbf{Missing}: energy, PUE, grid, idle, embodied}
    \end{column}
  \end{columns}
  \mlsyscredit{huggingface.co/bigscience/bloom @ 7f10a99, README.md, quoted exactly (source URL elided, long line wrapped; Jean Zay sentence omitted). Retrieved 19 Oct 2026.}
  \note{[1.5 min] A real model card: the README of BLOOM on Hugging Face, the document a training team publishes beside its weights. Two excerpts, quoted exactly. The metadata block, which tools read, gives emissions \val{card_g}, in grams: \val{card_t} tonnes, the dynamic consumption from the previous slide and \val{card_share} percent of the \val{b_total}-tonne life cycle (calcs.py). It names the source paper, the location, and the hardware. The prose section, which people read, still says ``Forthcoming'' for both carbon and electricity; the line omitted here says the Jean Zay supercomputer uses mostly nuclear energy. ENGAGE: one minute; which line of the record is missing? Step 2: energy at a stated boundary, PUE, grid intensity, idle power, and manufacturing. A reader comparing this card with another model's would compare the dynamic term with whatever the other team chose to report. WARN: a model card is the place to report carbon, but the field holds one number with no boundary. LINK: the synthesis writes the full record. FLEX: [OPTIONAL]; cutting it loses the model-card check.}
\end{frame}

% =============================================================================
% SECTION: OVER THE SYSTEM'S LIFE
% =============================================================================
\mlsysnote{[0.5 min] One training run is one term. A service runs for years. Which term dominates its lifetime cost and carbon? FLEX: [CORE].}
\mlsysdivider{Over three years, where do the dollars and the carbon go?}{Responsible Engineering, \S15.4.2, \S15.4.3, \S15.6}
\section{Lifetime Cost}

\begin{frame}{A busy ranking service keeps \val{gpus_0} GPUs busy around the clock}
  \centering
  \mlsyssteps[width=0.96\textwidth,height=3.1cm,keepaspectratio]{dlrm_strip_step}{4}\\[0.2cm]
  \begin{columns}[c]
    \begin{column}{0.30\textwidth}
      \centering\mlsysphoto[width=\textwidth,height=2.3cm,keepaspectratio]{datacenter-server-racks.jpg}
    \end{column}
    \begin{column}{0.36\textwidth}
      \small \val{gpus} GPUs is about \val{servers_w_words} \val{gpus_per_server}-GPU servers, busy every hour of the year.\\[0.1cm]
      \smalltag{Book scenario: \val{service_ms} ms of dedicated, unbatched GPU time per request.}
    \end{column}
    \begin{column}{0.28\textwidth}
      \uncover<4->{\begin{mlsyscard}{mlsysaccent}\small\textbf{One training cycle}\\ \$\val{cycle}\end{mlsyscard}}
    \end{column}
  \end{columns}
  \mlsyscredit{Book scenario: Responsible Engineering, \S15.4.2 (DLRM ranking service). Photo: Carl Lender, CC BY 2.0, Wikimedia Commons.}
  \note{[2.5 min] [CORE] Worked example 2, the book's DLRM scenario, which nobody measured. Step 1: \val{users_m} million users each get \val{items} ranked items a day: \val{inf_day_m} million inferences. Step 2: at \val{service_ms} ms of GPU time each, \val{inf_day_m} million times \val{service_s} seconds is \val{gpu_s_day_m} million GPU-seconds, divided by 3,600 gives \val{gpu_h_day} GPU-hours a day. Step 3: divided by 24 hours, \val{gpus} GPUs running continuously, about \val{servers_w_words} \val{gpus_per_server}-GPU servers (\val{servers} in calcs.py), racks like the ones in the photo. Step 4: at \$\val{inf_rate} per GPU-hour reserved, \val{gpus} times 24 times 365 times \val{inf_rate} is \$\val{serve_year} a year. One training cycle costs \$\val{cycle}: data preparation \$\val{cycle_prep}, hyperparameter search \$\val{cycle_search}, the final run \$\val{cycle_final}, at \$\val{train_rate} an hour on demand. WARN: \val{service_ms} ms of dedicated, unbatched GPU time per request is the chapter's assumption. LINK: Lecture 13's batching would shrink the GPU count; it would not change which term dominates. ENGAGE: before the next slide, guess the ratio of serving to training over three years. FLEX: [CORE].}
\end{frame}

\begin{frame}{Lifecycle TCO: serving is \val{cost_x} times training in cost, \val{carbon_x_dlrm} in carbon}
  \centering
  \mlsyssteps[width=0.90\textwidth,height=5.4cm,keepaspectratio]{tco_step}{2}
  \mlsyscredit{Book scenario: Responsible Engineering, table 10 and Napkin Math 15.4 (\val{gpu_w} W per GPU, \val{fleet_grid} kg/kWh, GPU only).}
  \note{[2 min] Step 1: total cost of ownership over three years. \Val{cycles} quarterly training cycles at \$\val{cycle}: \$\val{train_3y_k}K, \val{train_share} percent. Inference, \val{gpu_h_day} GPU-hours a day for \val{days} days at \$\val{inf_rate}: \$\val{inf_3y_m}M, \val{inf_share} percent. Operations, monitoring, retraining, and incident response at \$\val{ops_year_k}K a year: \$\val{ops_3y_k}K. Total \$\val{tco_m}M; serving against training is \val{cost_x} to 1. Step 2: carbon with Napkin Math 15.4. A \val{gpu_w} W GPU for one hour on a \val{fleet_grid} kg per kWh grid emits \val{gpu_kw} times \val{fleet_grid}, \val{kg_gpu_h} kg. Inference is \val{inf_gpu_h} GPU-hours (\val{gpu_h_day_1} a day for \val{days} days), \val{inf_t} tonnes; training \val{train_t} tonnes; \val{carbon_x_dlrm} times. Lifecycle total cost of ownership (TCO) is this sum over the system's whole life: training, inference, and operations. WARN: this is a GPU-only subtotal, no host, no PUE, no embodied term; the full boundary would raise both bars. Fallacy from \S15.6: measuring the environmental impact of training but not inference. FLEX: [CORE].}
\end{frame}

\begin{frame}{Estimate: cut serving \val{serve_cut_pct} percent, or halve all training?}
  \mlsysactivity{Estimate}{2.5 min}
  \begin{columns}[T]
    \begin{column}{0.40\textwidth}
      \begin{mlsyscard}{mlsysaccent}
        \textbf{Given (three years)}\\[0.1cm]
        Inference \$\val{inf_3y_m}M, \val{inf_t} t\\
        Training \$\val{train_3y_k}K, \val{train_t} t\\
        Total \$\val{tco_m}M
      \end{mlsyscard}
      \smalltag{Book scenario.}
    \end{column}
    \begin{column}{0.57\textwidth}
      \small
      \uncover<2->{Serving: $\val{serve_cut_frac} \times \val{inf_3y_m}\text{M} = \$\val{save_serve_k}$K;\quad $\val{serve_cut_frac} \times \val{inf_t} \approx \val{save_serve_t}$ t\\[0.1cm]
                   Training: $\val{train_cut_frac} \times \val{train_3y_k1}\text{K} = \$\val{save_train_k}$K;\quad $\val{train_cut_frac} \times \val{train_t} = \val{save_train_t}$ t\\[0.2cm]
      \centering\safeimg[width=\textwidth,height=2.9cm,keepaspectratio]{saving_step1.pdf}}
    \end{column}
  \end{columns}
  \mlsyscredit{Responsible Engineering, Checkpoint 15.4 and Summary.}
  \note{[2.5 min] ENGAGE: one minute; ask for the saving in dollars and the ratio. Step 2: a \val{serve_cut_pct} percent serving cut saves \$\val{save_serve_k}K, about \val{save_serve_pct} percent of the total, and about \val{save_serve_t} tonnes. Halving every training cycle saves \$\val{save_train_k}K and \val{save_train_t} tonnes, about \val{save_train_pct} percent of the total. The serving cut is worth \val{halvings} halvings of training. WARN: halving only the final run saves less still, \$\val{final_half} a cycle, because data preparation and the search remain; Checkpoint 15.4 asks exactly this. The ratio belongs to this scenario; a frontier model's training run can dominate instead. Epoch AI estimates frontier training cost grew about \val{frontier_growth} times a year, and the most expensive run in its data, Grok 4 in 2025, cost about \$\val{priciest_m}M (chapter figure 5). Measure which term dominates; do not assume it. FLEX: [CORE].}
\end{frame}

\begin{frame}{Real fleets agree, and an idle model still draws power}
  \begin{columns}[c]
    \begin{column}{0.47\textwidth}
      \centering
      \mlsyssteps[width=\textwidth,height=4.6cm,keepaspectratio]{fleet2_step}{2}
    \end{column}
    \begin{column}{0.50\textwidth}
      \uncover<3->{\small\textbf{BLOOM's public API, \val{api_days} days, no batching}\\[0.2cm]
      \begin{tabular}{@{}rl@{}}
        \textbf{\val{api_req}} & requests \\
        \textbf{\val{api_kwh} kWh} & on \val{api_gpus} A100 GPUs \\
        \textbf{\val{api_wh} Wh} & per request \smalltag{(derived)} \\
        \textbf{$\sim$\val{api_loaded_pct}\%} & to keep the model loaded, idle
      \end{tabular}}
    \end{column}
  \end{columns}
  \mlsyscredit{Patterson et al. 2022 (Google, 2019 to 2021); Wu et al., ``Sustainable AI,'' MLSys 2022, Fig. 3; Luccioni et al., JMLR 2023, \S4.4.}
  \note{[2 min] [OPTIONAL] Measured shares from two large operators back the scenario. Step 1: at Google, inference was about three fifths of ML energy from 2019 to 2021, and ML was \val{g_ml_lo} to \val{g_ml_hi} percent of total energy use. Step 2: Meta's AI power capacity split about \val{m_exp} to \val{m_train} to \val{m_inf} across experimentation, training, and inference. Step 3: one deployed model, measured. BLOOM's inference API on \val{api_gpus} A100 GPUs answered \val{api_req} requests in about \val{api_days} days, unbatched, using \val{api_kwh} kWh: \val{api_wh} Wh per request, derived by division. Even with almost no requests the instance drew about \val{api_idle_kwh} kWh every 10 minutes; the authors estimate about \val{api_loaded_pct} percent of the energy kept the model in memory while waiting. WARN: \val{api_wh} Wh is not comparable with Gemini's \val{gem_wh} Wh; the boundary, the hardware, and the load differ. LINK: an endpoint that sits idle between requests is the utilization problem Lecture 13's batching solves. FLEX: [OPTIONAL]; drop the left panel to save a minute.}
\end{frame}

\begin{frame}{True or false: newer servers always cut emissions}
  \mlsysactivity{Vote}{1.5 min}
  \centering
  {\small Replacing working servers with ones 10\% more efficient always lowers lifetime emissions.}\\[0.05cm]
  {\normalsize\alt<2->{\choice{True}{}\quad\rchoice{False}{}}{\choice{True}{}\quad\choice{False}{}}}\\[0.1cm]
  \mlsyssteps[width=0.92\textwidth,height=3.8cm,keepaspectratio]{embodied_step}{2}\\[0.02cm]
  {\footnotesize Derived: \val{gpu_w} W GPU; embodied \val{gpu_g_h} g per GPU-hour $+$ \val{server_g_h} g per server-hour $/$ \val{gpus_per_server} GPUs (assumed) $=$ \val{emb_g} g}
  \mlsyscredit{Derived from Luccioni et al. 2023 \S4.1 embodied factors, \val{gpus_per_server} GPUs per server (assumed), \val{gpu_w} W (Napkin Math 15.4). Gupta et al., ``Chasing Carbon,'' IEEE Micro 2022.}
  \note{[3 min] ENGAGE: vote by hands before step 2. Step 1 shows only the operational carbon of one GPU-hour: \val{gpu_w} W for an hour on a \val{fleet_grid} kg per kWh grid is \val{op_dirty_g} g; on France's \val{bloom_grid} grid, \val{op_clean_g} g. Step 2: add manufacturing. From BLOOM's factors, \val{gpu_g_h} g per GPU-hour plus a server's \val{server_g_h} g per hour shared by \val{gpus_per_server} GPUs, an assumption, gives \val{emb_g} g (calcs.py). On the dirty grid embodied is \val{emb_dirty_pct} percent; on the clean grid, \val{emb_clean_pct} percent. Answer: false. New hardware arrives with its manufacturing carbon already spent; on a clean grid a 10 percent cut in the operational term may never repay it (Gupta et al. 2022; chapter \S15.4.3). Measured check: for Meta's large-scale ML tasks, embodied was about \val{m_emb} percent of the footprint (Wu et al. 2022). Over the DLRM's three years, \val{inf_gpu_h} GPU-hours at \val{emb_g} g is about \val{dlrm_emb_t} tonnes embodied beside \val{inf_t0} operational. WARN: students treat electricity as the only carbon; on clean grids manufacturing decides. Name the vocabulary, for reading: scope 1, 2, and 3 accounting, life-cycle assessment, and carbon-aware scheduling, which moves flexible work to cleaner hours or regions. Land the segment: find the term that dominates before optimizing, and count manufacturing. FLEX: [OPTIONAL]; cutting it loses the embodied-carbon vote.}
\end{frame}

% =============================================================================
% SECTION: DUTIES HAVE MACHINE COSTS
% =============================================================================
\mlsysnote{[0.5 min] The system also owes duties to the people whose data it holds. Each duty becomes something the machine computes, stores, or staffs. FLEX: [CORE].}
\mlsysdivider{What does a legal duty become inside the machine?}{Responsible Engineering, \S15.5, \S15.5.3, \S15.3.3}
\section{Duties}

\begin{frame}{No breach, and still a \val{dpc_total} million euro fine}
  \begin{columns}[T]
    \begin{column}{0.42\textwidth}
      \centering
      \bignum{\texteuro \val{dpc_total}M}{Facebook \texteuro \val{dpc_fb}M + Instagram \texteuro \val{dpc_ig}M}\\[0.3cm]
      \raggedright{\small Finding: no valid legal basis for behavioral advertising; transparency and fairness breached. No data was stolen.}
    \end{column}
    \begin{column}{0.55\textwidth}
      \begin{mlsyscard}{midgray}
        \footnotesize\textbf{Data Protection Commission announces conclusion of two inquiries into Meta Ireland}\\[0.05cm]
        {\scriptsize\textcolor{midgray}{04th January 2023}}\\[0.1cm]
        \scriptsize ``Final decisions have now been made by the DPC in which it has fined Meta Ireland \texteuro 210 million (for breaches of the GDPR relating to its Facebook service), and \texteuro 180 million (for breaches in relation to its Instagram service).''
      \end{mlsyscard}
    \end{column}
  \end{columns}
  \vspace{0.25cm}
  \centering\uncover<2->{\normalsize For each use of each record: which legal basis allowed it,\\and where does the system record that?}
  \mlsyscredit{Data Protection Commission (Ireland), press release, 4 January 2023, dataprotection.ie/en/news-media/data-protection-commission-announces-conclusion-two-inquiries-meta-ireland, quoted exactly.}
  \note{[2.5 min] [CORE] The card quotes the regulator's own press release, its headline, date, and the paragraph that states the fines. Step 1: Ireland's Data Protection Commission fined Meta Ireland \val{dpc_fb} million euro for Facebook and \val{dpc_ig} million for Instagram, announced 4 January 2023 (the DPC's own release). Under GDPR, processing personal data is lawful only under one of six legal bases. Meta relied on ``contract'' for behavioral advertising; the European Data Protection Board ruled that basis invalid for that purpose, and the DPC also found transparency and fairness breaches. The Board directed the DPC to raise the fines it had proposed. Nothing was hacked. Step 2: the systems question. Security decides who can read data; governance must also prove why each use was lawful. ENGAGE: would encrypting every table have changed the finding? Expected: no; encryption controls access and says nothing about purpose. WARN: this lecture teaches capabilities to build and gives no legal advice; the law varies by jurisdiction and use, and the decision belongs with counsel. FLEX: [CORE].}
\end{frame}

\begin{frame}{Each obligation becomes a capability with a machine cost}
  \begin{columns}[c]
    \begin{column}{0.74\textwidth}
      \footnotesize
      \renewcommand{\arraystretch}{1.0}
      \begin{tabular}{@{}P{0.33\textwidth}P{0.31\textwidth}P{0.28\textwidth}@{}}
        \toprule
        \textbf{Obligation} & \textbf{Capability} & \textbf{Machine cost} \\
        \midrule
        Reasons for denial (ECOA; GDPR Art.\ 15) & Reason codes per decision & \leavevmode\uncover<3->{Explanation compute} \\
        Contest a decision (GDPR Art.\ 22) & Staffed human review & \leavevmode\uncover<3->{Reviewer hours} \\
        \leavevmode\uncover<2->{Erasure (GDPR)} & \leavevmode\uncover<2->{Lineage-driven deletion} & \leavevmode\uncover<3->{Retraining hours} \\
        \leavevmode\uncover<2->{Records (EU AI Act)} & \leavevmode\uncover<2->{Versions, provenance, logs, model card} & \leavevmode\uncover<3->{Log bytes} \\
        \textcolor{midgray}{Nondiscrimination} & \textcolor{midgray}{Per-group gates} & \textcolor{midgray}{Labeled data (Part 1)} \\
        \bottomrule
      \end{tabular}
    \end{column}
    \begin{column}{0.24\textwidth}
      \centering
      \uncover<2->{\mlsysphoto[width=\textwidth,height=2.9cm,keepaspectratio]{amazon-echo-dot-board.jpg}\\[0.05cm]
      {\scriptsize Smart-speaker main board: the wake-word detector runs here.}}
    \end{column}
  \end{columns}
  \mlsyscredit{Responsible Engineering, table 11 and \S15.5.3. Photo: Raimond Spekking (CC BY-SA 4.0), Wikimedia Commons.}
  \note{[2 min] Chapter table 11, built row by row. Step 1: the credit service. Under the Equal Credit Opportunity Act and its Regulation B, a denied US applicant is owed the specific principal reasons, and they must describe the factors actually used; GDPR Article 15 gives access to meaningful information about the logic. Article 22 lets people contest certain solely automated decisions. Step 2: the voice assistant. Erasure means finding every artifact derived from a recording; the EU AI Act, in force since August 1, 2024 with obligations phasing in, requires high-risk systems to keep versioned models, data provenance, and event logs. Step 3: the cost column. Every capability has a number: explanation compute, reviewer hours, retraining hours, log bytes. The next slides compute three of them. The gray row is Part 1's per-group release gate. Pitfall from \S15.6: documentation is not enforcement; a model card that says ``not validated for high-stakes decisions'' does nothing unless a deployment gate reads it. Lecture 14's promotion gate reads it; Lecture 12b set the acceptance gates it checks. WARN: not legal advice. FLEX: [CORE].}
\end{frame}

\begin{frame}{Estimate: what does a reason for every denial cost?}
  \mlsysactivity{Estimate}{3 min}
  \begin{columns}[T]
    \begin{column}{0.44\textwidth}
      \begin{mlsyscard}{mlsysaccent}
        \textbf{Given (credit service)}\\[0.1cm]
        \val{decisions_m} million decisions a day\\
        \val{deny_pct}\% denied\\
        Sampling explainer: \val{samples} evaluations\\
        Exact Shapley: \val{features} features
      \end{mlsyscard}
      \smalltag{Constructed example.}
    \end{column}
    \begin{column}{0.53\textwidth}
      \small
      \uncover<2->{Notices: $10^6 \times \val{deny_frac} = \val{notices}$ a day\\[0.15cm]}
      \uncover<3->{Sampled: $\val{notices} \times \val{samples} = \mathbf{\val{sampled_m}}$ \textbf{million}\\
                   \quad $= \val{sampled_x}\times$ the evaluations that made the decisions\\[0.15cm]}
      \uncover<4->{Exact: $\val{notices} \times 2^{\val{features}} \approx \mathbf{\val{exact_mant} \times 10^{\val{exact_exp}}}$}
    \end{column}
  \end{columns}
  \mlsyscredit{Responsible Engineering, \S15.3.3 (explanation cost) and \S15.5.3; denial share from the worked example's matrices (\val{credit_appr_n} of \val{credit_all_n} approved).}
  \note{[2.5 min] Worked example 3. ENGAGE: one minute; ask for evaluations a day and the ratio to the decisions. The credit service makes \val{decisions_m} million decisions a day; in the chapter's worked example \val{credit_appr_n} of \val{credit_all_n} applicants were approved, so \val{deny_pct} percent are denied, and each denial is owed reasons. Step 2: \val{notices} adverse-action notices a day. Step 3: a sampling explainer such as KernelSHAP evaluates the model on \val{samples} perturbed inputs per explanation: \val{sampled_m} million evaluations a day, \val{sampled_x} times the decision compute. Step 4: exact Shapley attribution evaluates every subset of \val{features} features, 2 to the \val{features}th, \val{exact_per}, so about \val{exact_mant} times 10 to the \val{exact_exp}th a day. WARN: the common error is to explain every decision, giving \val{all_decisions_b} billion; only denials need notices. FLEX: [CORE].}
\end{frame}

\begin{frame}{The explanation path can cost more than the model itself}
  \centering
  \mlsyssteps[width=0.94\textwidth,height=3.9cm,keepaspectratio]{explain_step}{3}\\[0.02cm]
  {\footnotesize Credit service (constructed): \val{deny_pct}\% of \val{decisions_m} million daily decisions are denials.}\\[0.15cm]
  \uncover<4->{\small Or choose a model that gives its own reasons, since a sparse linear model needs no explainer;\\TreeSHAP is exact for trees in polynomial time.}
  \note{[1 min] The estimate on one log axis. Step 1: the decisions, \val{decisions_m} million evaluations. Step 2: sampled explanations, \val{sampled_m} million. Step 3: exact attribution, \val{exact_mant} times 10 to the \val{exact_exp}th. A decision that takes one evaluation to make takes a million to explain exactly. Step 4: the alternative is a design choice made before training. An inherently interpretable model, such as a sparse linear model, yields reasons directly, at whatever accuracy that structure gives up; model-specific methods such as TreeSHAP exploit a tree's structure. WARN: an attribution is not a causal reason; Regulation B requires reasons that describe the factors actually used, so the explainer's fidelity must itself be tested. The chapter's interpretability spectrum (figure 4) places model families from interpretable to post-hoc. FLEX: [OPTIONAL]; fold step 4 into the estimate if short.}
\end{frame}

\begin{frame}{\val{review_h} reviewer-hours: \val{appeals} real reviews or \val{cigna_cases} rubber stamps}
  \centering
  {\small Appeals: $\val{appeal_pct}\% \times 10^6 = \val{appeals}$ a day.\quad $\val{appeals} \times \val{case_min}$ min $= \mathbf{\val{review_h}}$ \textbf{reviewer-hours a day} \smalltag{(\val{case_min} min assumed)}\\$\approx \val{reviewers}$ reviewers on \val{shift_h}-hour shifts.}\\[0.1cm]
  \mlsyssteps[width=0.94\textwidth,height=4.95cm,keepaspectratio]{review_zone_step}{2}
  \mlsyscredit{Responsible Engineering, \S15.5.3. Rucker, Miller, Armstrong, ProPublica and The Capitol Forum, March 25, 2023 (Cigna).}
  \note{[2 min] Contestability needs people. The curve is fixed hours divided by volume: as cases per reviewer-hour grow, time per case falls until review drops below any substantive minimum, into the rubber-stamp zone; the dashed line's position is illustrative. Step 1: if \val{appeal_pct} percent of \val{decisions_m} million daily decisions are appealed, the chapter's workload assumption, reviewers face \val{appeals} cases a day. At \val{case_min} minutes each, our assumption, that is \val{review_h} reviewer-hours a day, about \val{reviewers} full-time reviewers. Step 2: the same \val{review_h} hours, measured elsewhere. ProPublica reported that over \val{cigna_months} months Cigna's medical directors denied more than \val{cigna_cases} payment requests, spending on average \val{cigna_s} seconds on each. \val{cigna_cases} times \val{cigna_s} seconds is \val{cigna_h} hours. ENGAGE: what would make the \val{appeals} reviews real? Expected: tooling that shows each case's explanation, lineage, and model version, and a staffing budget sized to the volume. WARN: a human in the loop is not oversight by itself; the European Data Protection Board says required review must be substantive and must never become rubber-stamping. Land the segment: each duty is a number: explanation compute per denial, reviewer hours per appeal, a recorded legal basis per use. FLEX: [OPTIONAL]; cutting it loses the human-review budget.}
\end{frame}

% =============================================================================
% SECTION: GOVERNANCE AS A DATA SYSTEM
% =============================================================================
\mlsysnote{[0.5 min] Nearly every obligation asks the infrastructure two questions: which data made this model, and who touched it, under what authority. Add a third: once the data is in the weights, can it be removed? FLEX: [CORE].}
\mlsysdivider{Which data made this model, can we remove it, and can we prove it?}{Responsible Engineering, \S15.5.2, \S15.5.4, \S15.5.5, \S15.6}
\section{Governance}

\begin{frame}{An erasure request walks the lineage graph to the devices}
  \centering
  \mlsyssteps[width=0.94\textwidth,height=5.2cm,keepaspectratio]{lineage_step}{4}
  \mlsyscredit{Responsible Engineering, \S15.5 and \S15.5.4 (KWS pipeline). GDPR Art.\ 12(3), Regulation (EU) 2016/679.}
  \note{[2 min] [CORE] The voice assistant's pipeline as a lineage graph: artifacts are nodes, the steps that made them are edges. Step 1: a recording enters with its consent record. Step 2: it produces a spectrogram cache and feature embeddings. Step 3: those feed a training snapshot and checkpoint v12; a signed manifest of content hashes and augmentation seeds links the checkpoint to its exact training split. Step 4: deployment edges record which devices downloaded v12. An erasure request starts at the recording and follows every edge. Caches and embeddings are deleted, the snapshot is rebuilt, the weights need retraining, validated unlearning, or another control, and the devices can be updated. GDPR gives one month to respond, extendable by two more. LINK: Lecture 4 built this graph for reproducibility and Lecture 14 the model registry; governance adds only the query. WARN: finding an artifact does not by itself require deleting it; what is in scope is a legal decision (self-check 15.5, question 2). FLEX: [CORE].}
\end{frame}

\begin{frame}{True or false: deleted records leave the model}
  \mlsysactivity{Vote}{1.5 min}
  \begin{columns}[c]
    \begin{column}{0.55\textwidth}
      \centering
      {\small A user's records are deleted from the training database. The deployed model no longer contains them.}\\[0.2cm]
      {\Large \alt<2->{\choice{True}{}\quad\rchoice{False}{}}{\choice{True}{}\quad\choice{False}{}}}\\[0.2cm]
      \uncover<2->{{\fontsize{32}{36}\selectfont\bfseries\textcolor{mlsysaccent}{\val{carlini_mem}}}\\[0.1cm]
      {\footnotesize verbatim training texts extracted from GPT-2, out of \val{carlini_cand} candidates; \val{carlini_names} named individuals, \val{carlini_contact} with contact details}}
    \end{column}
    \begin{column}{0.42\textwidth}
      \centering
      \uncover<2->{\safeimg[width=\textwidth,height=4.7cm,keepaspectratio]{carlini2021-fig1.png}\\[0.03cm]
      {\scriptsize Black bars: the authors' redactions.}}
    \end{column}
  \end{columns}
  \mlsyscredit{Carlini et al., ``Extracting Training Data from Large Language Models,'' USENIX Security 2021; figure: arXiv 2012.07805v2, Figure 1, cropped.}
  \note{[1.5 min] [CORE] ENGAGE: vote by hands. The figure, after the vote, is the paper's Figure 1: a prefix fed to GPT-2 and the text it produced, a person's name, an email address, a phone and a fax number, and a physical address, all accurate, so the authors blacked them out. Step 2: false. Carlini and colleagues generated text from GPT-2, ranked \val{carlini_cand} candidate outputs, and confirmed \val{carlini_mem} were memorized verbatim from the training data; \val{carlini_names} contained individuals' names and \val{carlini_contact} contact information. Data trained into weights is still data. Membership inference attacks (Shokri et al. 2017) measure a weaker leak: whether a given record was in the training set. WARN: an attack that fails does not prove the data is gone; no empirical attack can prove a privacy bound. Fallacy from \S15.6: model weights are exempt from data governance. FLEX: [CORE].}
\end{frame}

\begin{frame}{Removing a record from the weights costs a retraining}
  \begin{columns}[c]
    \begin{column}{0.48\textwidth}
      \begin{mlsyscard}{mlsysaccent}
        \small\textbf{Retrain at GPT-3 scale}\\ \val{gpt3_mwh} MWh and \val{gpt3_t} t, every time
      \end{mlsyscard}
      \vspace{0.15cm}
      \begin{mlsyscard}{mlsysaccent}
        \small\textbf{DLRM: wait for the next cycle}\\ \$\val{cycle}, up to \val{cycle_days_r} days later
      \end{mlsyscard}
    \end{column}
    \begin{column}{0.50\textwidth}
      \centering
      \uncover<2->{{\small\textbf{Design for removal: SISA sharding}}\\[0.1cm]
      \safeimg[width=\textwidth,height=3.6cm,keepaspectratio]{sisa_step1.pdf}}
    \end{column}
  \end{columns}
  \mlsyscredit{Responsible Engineering, Napkin Math 15.5, \S15.4.2, \S15.5.5, \S15.6. Bourtoule et al., ``Machine Unlearning,'' IEEE S\&P 2021 (abstract).}
  \note{[2.5 min] Step 1: the reference method is to retrain without the record. At GPT-3 scale that re-spends \val{gpt3_mwh} MWh and \val{gpt3_t} tonnes for each request, which links erasure back to the first segment's carbon. For the DLRM the request can be folded into the next quarterly retraining, \$\val{cycle}, but the old model keeps serving for up to \val{cycle_days_r} days, and whether that delay is acceptable is a legal question. Step 2: machine unlearning designs training so removal is cheap. SISA splits the data into shards, trains a model per shard, and retrains only the shard that held the record: \val{sisa_purchase} times faster than retraining on Purchase, \val{sisa_svhn} on SVHN, \val{sisa_imagenet} on ImageNet, with a small accuracy loss. ENGAGE: why is the ImageNet gain the smallest? Expected: harder tasks need more data per model, so fewer, larger shards. WARN: verifying that an unlearning method removed a record's influence is an open problem (chapter \S15.6). FLEX: [CORE].}
\end{frame}

\begin{frame}{DP-SGD clips each example's gradient, then adds noise}
  \centering
  \mlsyssteps[width=0.94\textwidth,height=4.1cm,keepaspectratio]{dpsgd_mech_step}{4}\\[0.1cm]
  {\footnotesize Differential privacy: $\Pr[M(D) \in S] \le e^{\varepsilon}\,\Pr[M(D') \in S] + \delta$ for any datasets $D, D'$ that differ in one record.\\Smaller $\varepsilon$: less any one record can change the output.}
  \mlsyscredit{Responsible Engineering, \S15.5.2, after Abadi et al., ``Deep Learning with Differential Privacy,'' CCS 2016.}
  \note{[2.5 min] Differential privacy bounds what any one record can change about what a training run outputs: for any set of outputs S, adding or removing one person's record changes its probability by at most a factor of e to the epsilon, plus a small slack delta. Smaller epsilon is more private. At epsilon \val{eps}, an output can be up to e to the \val{eps}th, about \val{e_eps} times, more likely with the record than without it, which is why the accuracy slide calls \val{eps} a weak guarantee. DP-SGD changes one training step. Step 1: compute each example's gradient separately, before any batch average. Step 2: clip each to a maximum norm C, so no single example can pull the weights far. Step 3: add the clipped gradients and add Gaussian noise with standard deviation sigma times C, which hides any one example's contribution; then update. Step 4: a privacy accountant adds up the privacy spent at every step; training stops when the chosen epsilon is used up. ENGAGE: why must clipping come before the noise? Expected: the noise is scaled to C; without a bound on each example's influence, no finite noise hides it. WARN: differential privacy bounds influence under stated assumptions; it is not encryption. The definition is the chapter's footnote 16. FLEX: [CORE].}
\end{frame}

\begin{frame}[fragile]{Opacus clips and adds noise; per-example gradients cost memory}
  \begin{columns}[T]
    \begin{column}{0.54\textwidth}
      {\footnotesize\textbf{Opacus} README, ``Getting started''}
\vspace{-0.1cm}
\begin{semiverbatim}\fontsize{7.2}{8.2}\selectfont
privacy_engine = PrivacyEngine()
model, optimizer, data_loader =
  privacy_engine.make_private(
    module=model,
    optimizer=optimizer,
    data_loader=data_loader,
\textcolor{RedLine}{    noise_multiplier=1.1,}
\textcolor{OrangeLine}{    max_grad_norm=1.0,}
)
\end{semiverbatim}
      {\footnotesize\textbf{examples/mnist.py}, line 75}
\vspace{-0.1cm}
\begin{semiverbatim}\fontsize{7.2}{8.2}\selectfont
\textcolor{BlueLine}{epsilon = privacy_engine.accountant}
\textcolor{BlueLine}{  .get_epsilon(delta=args.delta)}
\end{semiverbatim}
    \end{column}
    \begin{column}{0.44\textwidth}
      \centering
      {\footnotesize\textcolor{OrangeLine}{\textbf{max\_grad\_norm}} $= C$, step 2\\ \textcolor{RedLine}{\textbf{noise\_multiplier}} $= \sigma$, step 3\\ \textcolor{BlueLine}{\textbf{accountant}}: step 4}\\[0.15cm]
      \safeimg[width=\textwidth,height=3.0cm,keepaspectratio]{dpsgd_step1.pdf}\\[0.1cm]
      {\footnotesize Dashed line: one H100, \val{h100_gb} GB.}
    \end{column}
  \end{columns}
  \mlsyscredit{github.com/pytorch/opacus @ v1.6.0: README.md lines 69 to 76 (one line wrapped) and examples/mnist.py line 75 (wrapped), quoted exactly. Memory: Responsible Engineering, \S15.5.2.}
  \note{[2 min] The real library call behind the diagram: Opacus, PyTorch's differential privacy library, version 1.6.0, quoted from its README and its MNIST example. make\_private wraps the model, the optimizer, and the data loader. max\_grad\_norm is the clipping bound C of step 2; noise\_multiplier is sigma of step 3, so the noise has standard deviation 1.1 times 1.0; the accountant of step 4 answers get\_epsilon for a chosen delta. To clip per example, the wrapped model attaches hooks that keep each example's gradient; that is the memory on the right. The \val{kws_params}-parameter KWS model's FP32 gradient is \val{kws_params} times \val{fp32_b} bytes, \val{kws_grad_kb} KB; a batch of \val{batch} is \val{kws_batch_mb} MB, easy. GPT-2 at \val{gpt2_params_b} billion parameters would need \val{gpt2_grad_gb} GB per example and \val{gpt2_batch_tb} TB per batch (\val{gpt2_batch_tb3} in calcs.py), beyond one \val{h100_gb} GB H100. Specialized implementations compute the clipping factors without keeping every gradient, or run small microbatches in sequence, at a cost in parallelism. ENGAGE: which term of memory grew? Expected: per-example state, which scales with the batch size as well as the model. FLEX: [OPTIONAL]; cutting it loses the real Opacus code.}
\end{frame}

\begin{frame}{Privacy at epsilon \val{eps}: over \val{cifar_gap} points from scratch, \val{im_gap_same} pre-trained}
  \centering
  \safeimg[width=0.94\textwidth,height=4.3cm,keepaspectratio]{dp_accuracy_step1.pdf}\\[0.05cm]
  {\footnotesize ImageNet: an NFNet-F3 pre-trained on JFT-4B, 4 billion images, then fine-tuned with DP-SGD.\\ The often-quoted \val{im_gap_best}-point gap is to the best non-private model of any kind.}
  \mlsyscredit{De et al., ``Unlocking High-Accuracy Differentially Private Image Classification through Scale,'' 2022 (arXiv 2204.13650v2): abstract, \S3, Table 6, Fig. 1; non-private best \val{im_best}\% from Yu et al. 2022.}
  \note{[1.5 min] Measured accuracy under differential privacy, two regimes. CIFAR-10, trained from scratch: De and colleagues reached \val{cifar_dp} percent at epsilon \val{eps}, \val{cifar_gain} points above the previous private best of \val{cifar_prev}, by tuning DP-SGD carefully on a large network; the same architecture, a Wide-ResNet 40-4, exceeds \val{cifar_np} percent without privacy, so privacy still costs over \val{cifar_gap} points (calcs.py). ImageNet: they did not train from scratch. They fine-tuned an NFNet-F3 already pre-trained, without privacy, on JFT-4B, a dataset of 4 billion labeled images, and reached \val{im_dp} percent at epsilon \val{eps}. The same pre-trained model fine-tuned without privacy reaches \val{im_same}, \val{im_gap_same} points higher (Table 6). The paper's headline, \val{im_gap_best} points, is the distance to \val{im_best} percent, the best non-private ImageNet result of any model at the time (Yu et al. 2022). ENGAGE: what closed most of the gap? Expected: pre-training on public data, which spends no privacy budget, then scale and careful tuning. WARN: do not quote \val{im_gap_best} points as the cost of privacy; it assumes a large public dataset, which a voice assistant or a credit model may not have. Epsilon \val{eps} is also a weak guarantee by the definition's standards; the same paper reports \val{im_eps05} percent at epsilon 0.5 when fine-tuning. Who chooses epsilon, and whether it is published, is open. FLEX: [OPTIONAL]; cutting it loses this measurement.}
\end{frame}

\begin{frame}{Raw audio stays on the device; model updates leave it}
  \centering
  \safeimg[width=0.94\textwidth,height=2.4cm,keepaspectratio]{kws_device.pdf}\\[0.1cm]
  \safeimg[width=0.94\textwidth,height=3.4cm,keepaspectratio]{kws_bytes_step1.pdf}
  \mlsyscredit{Responsible Engineering, \S15.5.2 (KWS, \val{kws_params_k}K parameters). Audio rate assumed: \val{audio_khz} kHz, \val{audio_bits}-bit; update = \val{kws_params_k}K FP32 parameters. Zhu et al., ``Deep Leakage from Gradients,'' NeurIPS 2019.}
  \note{[1.5 min] Privacy by architecture. The voice assistant's audio streams through a small circular buffer in on-chip SRAM, overwritten continuously, and may be saved or sent only after the on-device wake word fires. To still learn from many homes, federated learning trains on paired phones or hubs and sends model updates instead of audio. One hour of audio at \val{audio_khz} kHz and \val{audio_bits} bits, an assumption, is \val{audio_bps} bytes a second, \val{audio_mb_h} MB; one update of \val{kws_params} FP32 parameters is \val{kws_grad_kb} KB, \val{audio_x} times smaller. WARN: updates can still leak training examples (Zhu et al. 2019), so secure aggregation hides individual updates from the coordinator and differential privacy bounds what the final model reveals. False wakes release audio too; they must be tested. FLEX: [OPTIONAL]; cut first if short.}
\end{frame}

\begin{frame}{Estimate: how large is the credit service's audit log?}
  \mlsysactivity{Estimate}{2.5 min}
  \begin{columns}[T]
    \begin{column}{0.52\textwidth}
      \begin{mlsyscard}{mlsysaccent}
        \textbf{Given}\\[0.1cm]
        \val{decisions_m} million decisions a day\\
        One record per decision: \val{rec_b} B \smalltag{(assumed)}\\
        Keep \val{ret_months} months, about \val{ret_days} days
      \end{mlsyscard}
      \vspace{0.1cm}\small
      \uncover<2->{$10^6 \times \val{rec_b} \text{ B} = \val{log_mb_day}$ MB a day\\[0.05cm]}
      \uncover<3->{$\val{log_mb_day} \text{ MB} \times \val{ret_days} = \mathbf{\val{log_gb}}$ \textbf{GB}}
    \end{column}
    \begin{column}{0.45\textwidth}
      \centering
      \uncover<3->{\safeimg[width=\textwidth,height=3.6cm,keepaspectratio]{book_margin_004.pdf}\\[0.05cm]
      {\scriptsize Log volume $=$ events $\times$ bytes per event $\times$ retention}}
    \end{column}
  \end{columns}
  \mlsyscredit{Responsible Engineering, \S15.5.5 and table 12; Regulation B, 12 CFR 1002.12(b)(1) (\val{ret_months}-month retention).}
  \note{[2.5 min] Worked example 4. A decision manifest records what an auditor needs to reconstruct one decision: the model version hash, the threshold, the score or reason codes, and hashes of the input features. ENGAGE: one minute; ask for the total. Step 2: \val{decisions_m} million records of \val{rec_b} bytes, our assumption, is \val{log_mb_day} MB a day. Step 3: Regulation B requires US lenders to keep consumer credit application records for \val{ret_months} months, about \val{ret_days} days, so \val{log_gb} GB. The book's figure shows that the log grows with time while the dataset stays flat, and crosses the storage budget unless bytes per event stay small. WARN: a hash identifies an input but cannot reconstruct it; replay needs the source retained under its own rules. FLEX: [CORE].}
\end{frame}

\begin{frame}{Bytes per event is the term the design controls}
  \centering
  \mlsyssteps[width=0.94\textwidth,height=4.6cm,keepaspectratio]{audit_step}{2}\\[0.15cm]
  \uncover<2->{\small Log metadata and hashes; keep payloads only where a duty requires them.}
  \mlsyscredit{Responsible Engineering, \S15.5.5. DLRM input sized as in MLPerf DLRM on Criteo Terabyte: \val{dense} dense and \val{sparse} categorical features, \val{emb_dim}-dimensional FP32 embeddings (constructed).}
  \note{[1.5 min] Step 1: the credit service, \val{log_gb} GB. Step 2: the DLRM at \val{dlrm_req_m} million requests a day. Compact \val{rec_b}-byte manifests for a year: \val{dlrm_req_m} million times \val{rec_b} times \val{dlrm_days} is \val{dlrm_manifest_tb} TB. Logging the full input instead, \val{dense} dense features and \val{sparse} embeddings of \val{emb_dim} FP32 values, \val{input_b} bytes, gives \val{dlrm_input_tb} TB, \val{payload_x} times more. Event count is set by the traffic; retention is set by law; bytes per event is the term engineers choose. Edge tier, for reading: the KWS device logs state changes to a volatile ring buffer and uploads while charging, and sequence numbers expose gaps; a gap cannot prove no incident occurred. Land the segment: lineage finds the data, weights can still hold it, removal costs a retraining or a design for unlearning, differential privacy pays in memory and accuracy, and the audit log stays bounded only if it records metadata. LINK: Lecture 14b owns incident response that reads these logs. FLEX: [OPTIONAL]; cutting it loses the bytes-per-event lever.}
\end{frame}

% =============================================================================
% CLOSE
% =============================================================================
\begin{frame}{The record that checks a carbon number also proves a duty}
  \begin{columns}[c]
    \begin{column}{0.40\textwidth}
      \centering
      {\fontsize{22}{26}\selectfont\bfseries\textcolor{midgray}{\val{nas_est_t} t}}\\[0.02cm]
      {\small estimated in 2019}\\[0.12cm]
      {\fontsize{22}{26}\selectfont\bfseries\textcolor{mlsysaccent}{\val{nas_meas_t} t}}\\[0.02cm]
      {\small measured in 2021: \val{nas_err}$\times$ lower}
    \end{column}
    \begin{column}{0.58\textwidth}
      \small
      \uncover<2->{\textbf{One run's record}\\[0.02cm]
      \renewcommand{\arraystretch}{1.0}
      \begin{tabular}{@{}lP{0.70\textwidth}@{}}
        \toprule
        Energy & kWh at a stated boundary \\
        PUE & of the datacenter used \\
        Hardware & processor and measured power \\
        Grid & intensity, and on which basis \\
        Lineage & proxy task, model size, checkpoints \\
        \bottomrule
      \end{tabular}}\\[0.15cm]
      \uncover<3->{\textcolor{mlsysaccent}{\textbf{Lineage answers the erasure query;\\the manifest answers the audit.}}}
    \end{column}
  \end{columns}
  \vspace{0.1cm}
  \uncover<2->{\safeimg[width=\textwidth,keepaspectratio]{patterson2021-appD-line.png}\\[0.02cm]
  {\footnotesize\textcolor{midgray}{days $\times$ hours $\times$ chips $\times$ watts $\times$ PUE $=$ \val{nas_kwh} kWh;\quad $\times$ \val{grid_g} kg/kWh $=$ \val{nas_calc_t} t}}}
  \mlsyscredit{Patterson et al., ``Carbon Emissions and Large Neural Network Training,'' arXiv 2104.10350v3, Appendix D, p.\ 22, the record as printed, one line cropped from the PDF.}
  \note{[2.5 min] Back to the hook. ENGAGE: what record, published in 2019, would have settled the estimate? Take two answers. Step 2: read the paper's own line at the bottom: \val{nas_days} days, 24 hours, \val{nas_chips} TPU v2 chips, \val{nas_w} W measured, PUE \val{pue_g}, so \val{nas_kwh} kWh, times \val{grid_g} kg per kWh, \val{nas_calc_t} tonnes (calcs.py). Every factor of the law is printed, so anyone can check it. The record: the measured energy and its boundary, the processor and its measured power, the datacenter's PUE, the grid intensity and whether it was counted market-based or location-based, and the lineage of the run: the proxy task, the model size, the checkpoints. Those are the three facts that separated \val{nas_est_t} tonnes from \val{nas_meas_t}: the proxy task, the TPU v2, and the PUE. The same record showed the search paid back \val{payback} to 1 (\val{payback_1} in calcs.py); Meena, built on the architecture it found, saved \val{meena_t} tonnes. Step 3: the same record, read from the governance side. Lineage from data to checkpoint to device is what an erasure request queries; a compact manifest per run or per decision is what an auditor reads. The same record (energy at a stated boundary, hardware, PUE, grid, lineage) settles the carbon question and answers an erasure request or an audit. FLEX: [CORE].}
\end{frame}

\begin{frame}{Write the ledger for one service}
  \mlsysactivity{Retrieval}{2 min}
  \small
  \renewcommand{\arraystretch}{1.15}
  \begin{tabular}{@{}P{0.30\textwidth}P{0.66\textwidth}@{}}
    \toprule
    \textbf{Line} & \leavevmode\uncover<2->{\textbf{DLRM or credit service}} \\
    \midrule
    Carbon a year, largest term & \leavevmode\uncover<2->{DLRM: \val{inf_t_year} t modeled (\val{inf_t} t in \val{years} years), inference} \\
    Erasure: path and cost & \leavevmode\uncover<2->{lineage to checkpoint and devices; next cycle or SISA} \\
    Audit record and size & \leavevmode\uncover<2->{manifests: \val{log_gb} GB (credit); payloads \val{payload_x}$\times$ more} \\
    Review staffing & \leavevmode\uncover<2->{\val{appeals} appeals, \val{review_h} reviewer-hours a day} \\
    \bottomrule
  \end{tabular}
  \note{[2 min] From memory, students write one line per row for either service; two minutes. Step 2: the answers. DLRM carbon: \val{inf_t} tonnes over three years of modeled accelerator carbon, about \val{inf_t_year} a year, almost all inference; embodied adds about \val{dlrm_emb_t_year} a year at the derived \val{emb_g} g per GPU-hour. Erasure: walk the lineage to the checkpoint and the devices; fold the request into the next quarterly cycle or design with SISA. Audit: compact manifests; the credit service keeps \val{log_gb} GB over \val{ret_months} months, and full DLRM payloads would be \val{payload_x} times the manifests. Review: \val{appeals} appeals a day at \val{case_min} minutes is \val{review_h} reviewer-hours. Ask which line students could not fill in; that is the one to reread. FLEX: [CORE].}
\end{frame}

\begin{frame}{Five things to carry out of this lecture}
  \thumbline{thumb_factors_step4.pdf}{Carbon is energy $\times$ PUE $\times$ grid intensity; model, processor, datacenter, and grid each move it by factors (\val{ff_grid}$\times$ on one job).}\\[0.2cm]
  \thumbline{thumb_bloom_life_step3.pdf}{A training-energy number leaves out idle power and manufacturing, half of BLOOM's \val{b_total} t.}\\[0.2cm]
  \thumbline{thumb_tco_step2.pdf}{For a busy service, serving dominates; a \val{serve_cut_pct}\% serving cut is worth \val{halvings} halvings of training.}\\[0.2cm]
  \thumbline{thumb_explain_step3.pdf}{Every duty has a machine cost: \val{sampled_x}$\times$ compute for reasons, \val{review_h} reviewer-hours for appeals.}\\[0.2cm]
  \thumbline{thumb_lineage_step4.pdf}{Lineage finds data, weights can hold it, privacy costs memory and accuracy, logs stay bounded by recording metadata.}
  \note{[1.5 min] Read the five lines. The first two belong to segment 1, the third to segment 2, the fourth to segment 3, the fifth to segment 4. Ask students to star the line they found least obvious and to check that number in the chapter tonight. FLEX: [CORE].}
\end{frame}

\begin{frame}{Questions with no single answer}
  \normalsize
  \begin{itemize}\setlength{\itemsep}{0.35cm}
    \item Should a service report carbon per request beside its latency? On which boundary: market or location, with or without manufacturing?
    \item An erasure request arrives for a model that cost \val{gpt3_mwh} MWh to train. Retrain, unlearn, or document an exception? Who decides?
    \item Who should choose epsilon for a voice assistant, and should the choice be published?
  \end{itemize}
  \note{[1 min] No single answer. Ask students to write one sentence on any of the three. The first returns to the Gemini and BLOOM boundaries; the second to the cost of retraining; the third to the accuracy and memory cost of privacy. FLEX: [OPTIONAL].}
\end{frame}

\begin{frame}[fragile]{Read, and measure your own run}
  \begin{columns}[T]
    \begin{column}{0.50\textwidth}
      \footnotesize
      \begin{itemize}\setlength{\itemsep}{0.04cm}
        \item \textbf{Read:} Responsible Engineering, \S15.4, \S15.5.2 to \S15.5.5, \S15.6; Napkin Math 15.4 and 15.5; Checkpoint 15.4; tables 11 and 12.
        \item Patterson et al. 2021, carbon of large training runs.
        \item Luccioni et al., JMLR 2023, BLOOM's life cycle.
        \item Elsworth et al. 2025, energy per Gemini prompt.
        \item Carlini et al. 2021; Dwork and Roth 2014 (privacy).
      \end{itemize}
    \end{column}
    \begin{column}{0.48\textwidth}
      {\footnotesize\textbf{CodeCarbon 3.3.1}, 30 s of matrix multiplies on a laptop}
\vspace{-0.1cm}
\begin{semiverbatim}\fontsize{7}{8.2}\selectfont
\textcolor{RedLine}{[codecarbon WARNING] No CPU tracking}
\textcolor{RedLine}{ mode found. Falling back on CPU load mode.}
duration: 30.075745915994048
cpu_power: 13.655516686012499
ram_power: 6.0
energy_consumed: 0.00018553642733992439
emissions: 6.855070041856389e-05
country_iso_code: USA
pue: 1.0
\end{semiverbatim}
      {\footnotesize \val{cc_wh} Wh and \val{cc_mg} mg: kWh and kg, at an estimated boundary.}
    \end{column}
  \end{columns}
  \mlsyscredit{CodeCarbon (MIT license), run here 19 Oct 2026; fields from its emissions.csv and its log, quoted exactly (log prefix time stamp removed, line wrapped).}
  \note{[1 min] [OPTIONAL] The chapter first, with its self-checks for sections 15.4 and 15.5; then the papers behind today's numbers. On the right, real output of an energy tracker, run for this lecture: pip install codecarbon, then OfflineEmissionsTracker(country\_iso\_code="USA") around 30 seconds of 2,048 by 2,048 float32 matrix multiplies on an Apple M5 Max laptop; the lines are fields of the emissions.csv it writes and one line of its log. Read the boundary before the number. The log says it found no power sensor and fell back on CPU load, so the CPU energy is an estimate; RAM is a fixed 6 W; PUE is 1.0, the tracker's default. Total \val{cc_wh} Wh, \val{cc_mg} milligrams, \val{cc_intensity} kg per kWh for the US grid it was told to assume (calcs.py). Assign it: wrap one training run of your own in the tracker and write down every assumption its output makes. FLEX: [OPTIONAL].}
\end{frame}

\begin{frame}[c]
  \centering
  {\small\textcolor{lightgray}{NEXT LECTURE \textbar{} CONCLUSION}}\\[0.4cm]
  \mlsyskitfig[width=0.30\textwidth,height=3.0cm,keepaspectratio]{stack}\\[0.4cm]
  {\LARGE\bfseries What does mastering the full stack let you do?}
  \note{[1 min] Every constraint so far, from latency to fairness to carbon, was taught one layer at a time. Lecture 16 follows one request across every layer and asks where the costly failures arise: where layers interact, when each component meets its own target and the system still misses its goal. Ask students to bring one number from each part of the course. FLEX: [CORE].}
\end{frame}

\end{document}
